% Sections 18.1 to 18.12, translated from sv/chapters/18/theory.tex.

\section{The standard library \code{math.h}}
\label{c18:sec:mathh}
\file{math.h} is the part of C's standard library that contains mathematical functions. The
functions work with floating-point numbers of type \code{double}.

\begin{ccode}
#include <math.h>
#include <stdio.h>
\end{ccode}

\subsection{Development environment}
\label{c18:sec:miljo}
During the lesson you use OnlineGDB, a \code{gcc} compiler that runs directly in the web browser:
\url{https://www.onlinegdb.com/online_c_compiler}. No installation is needed, and no account.
\begin{enumerate}
  \item Check that the language is set to \textbf{C} in the list at the top right.
  \item Write the code in the editor.
  \item Press \textbf{Run}. The program's output appears in the box below the editor.
\end{enumerate}

\begin{aside}
\note[Note] The code is not saved automatically. Copy it to a file of your own before you close
the tab.
\end{aside}

Programs that print sample values as CSV (\secref{c18:sec:csv}) are run in the usual way. Select
the output in the output box, copy it and paste it into a spreadsheet program to plot the signal.

\subsection{Compiling on your own computer (optional)}
\label{c18:sec:egendator}
If you would rather run the code locally, you can set up \textbf{WSL} (Windows Subsystem for
Linux) with \code{gcc}. Nothing in the chapter requires it -- OnlineGDB is enough all the way --
but the environment is used in later courses, so you can set it up now if you like. WSL runs a
Linux distribution, here Ubuntu, in a terminal directly in Windows, without a virtual machine.
The installation requires administrator rights and a restart.

Open \textbf{Windows PowerShell} as administrator and run:
\begin{shell}
wsl --install
wsl --set-default-version 2
\end{shell}
Restart the computer if the system asks you to. The command \code{wsl --status} should then show
\code{Default Version: 2}. If Ubuntu was not installed, the latest LTS version is in the Microsoft
Store. WSL requires Windows~11, or Windows~10 version 2004 (build 19041) or later, and
virtualisation enabled in the computer's UEFI/BIOS. Without administrator rights, WSL cannot be
installed; use OnlineGDB instead.

Then start \textbf{Ubuntu} from the Windows search box. On the first start you choose a user name
and a password; the password is not shown while you type it, which is normal. Update the system
and install the development tools, including \code{gcc} and \code{make}:
\begin{shell}
sudo apt -y update
sudo apt -y upgrade
sudo apt -y install build-essential
gcc --version
\end{shell}
The Windows files are reached from Ubuntu under \file{/mnt/c/}, but work is faster if the files
are in the Linux home directory, \code{~}. \textbf{VS Code} with the \textbf{WSL} extension opens
the files directly in the Ubuntu environment; \code{code .} in the terminal opens the current
directory.

\subsection{Linking with the maths library}
\label{c18:sec:lankning}
The functions in \file{math.h} and \file{complex.h} are in a separate maths library, \code{libm}.
On your own computer the program must therefore be linked with the flag \code{-lm}, which is
placed \textbf{last} on the command line, after the source files:
\begin{shell}
gcc main.c -o main -lm
./main
\end{shell}
Without \code{-lm} the program compiles, but linking fails with error messages such as
\code{undefined reference to 'sqrt'}. If you get that error in OnlineGDB, add \code{-lm} under
\textbf{Extra Compiler Flags} in the settings (the gear icon).

\subsection{Constants}
\label{c18:sec:konstanter}
\code{M_PI} ($\pi$) and \code{M_E} ($e$) are in \file{math.h} on most systems, but they are not
part of the C standard and are therefore missing in strict standard mode (\code{-std=c11}).
OnlineGDB uses \code{gcc}'s default mode, so \code{M_PI} works there, as it does with
\code{-std=gnu11}. For the code to work everywhere, the constant can be defined directly in the
program:
\begin{ccode}
#define PI 3.14159265358979323846
\end{ccode}

\section{Floating-point numbers in C: \code{float} and \code{double}}
\label{c18:sec:flyttal}
The integer types you have already used, such as \code{int}, \code{uint8_t} and \code{uint16_t},
can only store whole numbers: \code{7 / 2} becomes \code{3}, not $3.5$. The mathematics in this
course requires decimal numbers, and for those \textbf{floating-point numbers} are used.

\textbf{Why are floating-point numbers avoided in microcontroller code?} An AVR microcontroller has
no floating-point unit (FPU). Floating-point calculations must therefore be done in software,
which is slow and takes a lot of memory. On AVR, \code{double} is moreover often only 32~bits,
the same precision as \code{float}. An ordinary computer has an FPU and calculates with
floating-point numbers at no extra cost -- which is why they are used freely in this chapter.

\subsection{Types}
\label{c18:sec:typer}
\begin{booktable}{@{}p{4.5em}p{5em}p{12.5em}L@{}}
  \thead{Type} & \thead{Size} & \thead{Approximate precision} & \thead{Use} \\
  \midrule
  \code{float} & 32 bits & about 7 significant figures & When memory is limited \\
  \code{double} & 64 bits & about 15 significant figures & \textbf{The default in this course} \\
\end{booktable}

Every function in \file{math.h} takes and returns \code{double}. So use \code{double} throughout.

\subsection{Literals}
\label{c18:sec:litteraler}
How a number is written decides which type it gets:
\begin{ccode}
int    a = 2;      // Integer.
double b = 2.0;    // Floating point (double).
float  c = 2.0f;   // Floating point (float).
\end{ccode}

\begin{aside}
\note[Note] If two integers are divided, integer division is performed, whatever type the result
is stored in:
\begin{ccode}
double x = 1 / 2;      // 0.0 - integer division!
double y = 1.0 / 2.0;  // 0.5
\end{ccode}
\end{aside}

\subsection{Printing}
\label{c18:sec:utskrift}
\begin{narrowtable}{@{}ll@{}}
  \thead{Type} & \thead{Format specifier} \\
  \midrule
  \code{int}, \code{uint8_t}, \code{uint16_t} & \code{\%d} \\
  \code{float}, \code{double} & \code{\%f}, for example \code{\%.3f} for three decimals \\
\end{narrowtable}

\code{printf()} automatically converts \code{float} to \code{double}, so \code{\%f} works for
both. Never use \code{\%d} for a floating-point number -- the output is then meaningless.

\subsection{Converting between integers and floating-point numbers}
\label{c18:sec:omvandling}
Integers are converted to floating-point numbers automatically. In the other direction the number
is \textbf{truncated}, not rounded:
\begin{ccode}
const double d = 7;           // 7.0
const int    n = 3.9;         // 3, not 4!
const int    m = round(3.9);  // 4
\end{ccode}

\subsection{Precision}
\label{c18:sec:precision}
Floating-point numbers are stored in binary and cannot represent every decimal number exactly.
The expression \code{0.1 + 0.2} gives \code{0.30000000000000004}, not exactly \code{0.3}.

Never compare floating-point numbers with \code{==}, therefore. Check instead that the difference
is small enough:
\begin{ccode}
if (fabs(a - b) < 1e-9) { /* The numbers are considered equal. */ }
\end{ccode}

\section{Basic functions}
\label{c18:sec:funktioner}
\begin{booktable}{@{}>{\raggedright\arraybackslash}p{8.5em}p{7.5em}L@{}}
  \thead{Mathematics} & \thead{C function} & \thead{Comment} \\
  \midrule
  $\sqrt{x}$ & \code{sqrt(x)} & Square root, $x \geq 0$ \\
  $x^y$ & \code{pow(x, y)} & Power \\
  $e^x$ & \code{exp(x)} & Exponential function \\
  $\ln x$ & \code{log(x)} & Natural logarithm, $x > 0$ \\
  $\lg x$ & \code{log10(x)} & Base-10 logarithm, $x > 0$ \\
  $\abs{x}$ & \code{fabs(x)} & Absolute value of a floating-point number \\
  $\sqrt{x^2 + y^2}$ & \code{hypot(x, y)} & Magnitude of a vector \\
  Remainder of a division & \code{fmod(x, y)} & Like \code{\%}, but for floating-point numbers \\
  Rounding & \code{round(x)} & \code{floor(x)} down, \code{ceil(x)} up \\
\end{booktable}

\begin{aside}
\note[Note] \code{\%} works only for integers. For floating-point numbers, \code{fmod()} is used.
\end{aside}

\begin{ccode}
const double r = sqrt(pow(3.0, 2.0) + pow(4.0, 2.0)); // r = 5.0
printf("r = %.2f\n", r);
\end{ccode}

\section{Trigonometric functions}
\label{c18:sec:trig}
\begin{narrowtable}{@{}lll@{}}
  \thead{Mathematics} & \thead{C function} & \thead{Return value} \\
  \midrule
  $\sin(x)$ & \code{sin(x)} & -- \\
  $\cos(x)$ & \code{cos(x)} & -- \\
  $\tan(x)$ & \code{tan(x)} & -- \\
  $\arcsin(x)$ & \code{asin(x)} & $[-\pi/2,\ \pi/2]$ \\
  $\arccos(x)$ & \code{acos(x)} & $[0,\ \pi]$ \\
  $\arctan(x)$ & \code{atan(x)} & $(-\pi/2,\ \pi/2)$ \\
  $\arctan(y/x)$ & \code{atan2(y, x)} & $(-\pi,\ \pi]$ \\
\end{narrowtable}

\textbf{All the functions work in radians, never in degrees.} \code{sin(30)} is read as
$\sin(30\,\text{rad})$ and gives $-0.988$, not $0.5$.

\subsection{Converting between degrees and radians}
\label{c18:sec:radianer}
\[
  \text{rad} = \text{degrees} \cdot \frac{\pi}{180}, \qquad
  \text{degrees} = \text{rad} \cdot \frac{180}{\pi}
\]
\begin{ccode}
const double deg = 30.0;
const double rad = deg * PI / 180.0;
printf("sin(30 degrees) = %.3f\n", sin(rad)); // 0.500
\end{ccode}

\subsection{\code{atan2} instead of \code{atan}}
\label{c18:sec:atan2}
\code{atan(y / x)} cannot tell quadrant 1 from quadrant 3, or quadrant 2 from quadrant 4, because
the signs disappear in the division. \code{atan2(y, x)} takes $y$ and $x$ separately and therefore
gives the right angle in all four quadrants. Always use \code{atan2()} when converting from
rectangular to polar form (\kapref{15}).

\section{Powers, logarithms and decibels}
\label{c18:sec:decibel}
From \kapref{10}, for voltage and power respectively:
\[
  A_{\text{dB}} = 20\log_{10}\!\left(\frac{U_{\text{out}}}{U_{\text{in}}}\right), \qquad
  A_{\text{dB}} = 10\log_{10}\!\left(\frac{P_{\text{out}}}{P_{\text{in}}}\right)
\]
\begin{ccode}
const double gain_db = 20.0 * log10(u_out / u_in);
\end{ccode}
The way back from decibels to a ratio:
\[
  \frac{U_{\text{out}}}{U_{\text{in}}} = 10^{A_{\text{dB}}/20}
\]
\begin{ccode}
const double ratio = pow(10.0, gain_db / 20.0);
\end{ccode}
Logarithms to other bases are calculated with the change-of-base rule
$\log_b x = \frac{\ln x}{\ln b}$:
\begin{ccode}
const double log2_x = log(x) / log(2.0);
\end{ccode}

\section{Numerical differentiation}
\label{c18:sec:derivering}
The derivative is defined as the limit
\[
  f'(x) = \lim_{h \to 0} \frac{f(x + h) - f(x)}{h}.
\]
A computer cannot let $h \to 0$, but a small $h$ gives an approximate value. The \textbf{central
difference} gives better accuracy than the formula above, because the errors on either side of
$x$ cancel each other out:
\[
  f'(x) \approx \frac{f(x + h) - f(x - h)}{2h}
\]
\begin{ccode}
double derivative(double (*f)(double), const double x, const double h)
{
    return (f(x + h) - f(x - h)) / (2.0 * h);
}
\end{ccode}
The call is made with a function pointer, that is, the name of the function to be
differentiated:
\begin{ccode}
const double slope = derivative(square, 3.0, 1e-6);
\end{ccode}

\subsection{Choosing \texorpdfstring{$h$}{h}}
\label{c18:sec:valh}
\begin{itemize}
  \item An $h$ that is too \textbf{large} gives a truncation error -- the approximation lies too
    far from the limit.
  \item An $h$ that is too \textbf{small} gives a rounding error -- $f(x + h)$ and $f(x - h)$
    become so alike that the significant figures vanish in the subtraction.
\end{itemize}
For \code{double}, $h \approx 10^{-6}$ is a good compromise.

\section{Numerical integration}
\label{c18:sec:integrering}
The definite integral corresponds to the area under the curve (\kapref{14}). If the interval
$[a, b]$ is divided into $n$ equal parts of width $h = \frac{b - a}{n}$, the area can be
approximated with the \textbf{trapezoidal rule}:
\[
  \int_a^b f(x)\,dx \approx h\left[\frac{f(a) + f(b)}{2} + \sum_{k=1}^{n-1} f(a + kh)\right]
\]
Each subinterval is thus approximated by a trapezium instead of a rectangle, which follows the
curve considerably better. The accuracy increases with the number of subintervals $n$.
\begin{ccode}
double integral(double (*f)(double), const double a, const double b, const int n)
{
    const double h = (b - a) / n;
    double sum     = (f(a) + f(b)) / 2.0;

    for (int k = 1; k < n; ++k)
    {
        sum += f(a + k * h);
    }
    return sum * h;
}
\end{ccode}
For a first-degree polynomial the trapezoidal rule gives an \textbf{exact} result, since the curve
is then a straight line. For other functions there is an error, which decreases as $n$ increases.

\section{Complex numbers with \code{complex.h}}
\label{c18:sec:complex}
Complex numbers (Chapters~\ref{c15:ch}--\ref{c17:ch}) are handled by \file{complex.h}, which is
used together with \file{math.h}:
\begin{ccode}
#include <complex.h>
#include <math.h>
\end{ccode}
The type is called \code{double complex} and the imaginary unit is called \code{I}:
\begin{ccode}
const double complex z = 3.0 + 4.0 * I;
\end{ccode}

\begin{aside}
\note[Note] In electrical engineering the imaginary unit is written $j$, so that it is not
confused with the current $i$. In C it is always called \code{I}.
\end{aside}

\subsection{Functions}
\label{c18:sec:cfunktioner}
\begin{booktable}{@{}p{7em}p{7.5em}L@{}}
  \thead{Mathematics} & \thead{C function} & \thead{Description} \\
  \midrule
  $\text{Re}(z)$ & \code{creal(z)} & The real part \\
  $\text{Im}(z)$ & \code{cimag(z)} & The imaginary part \\
  $\abs{z}$ & \code{cabs(z)} & The magnitude \\
  $\arg(z)$ & \code{carg(z)} & The argument in radians, $(-\pi,\ \pi]$ \\
  $\bar{z}$ & \code{conj(z)} & The conjugate \\
  $e^z$ & \code{cexp(z)} & The exponential function \\
  $\sqrt{z}$ & \code{csqrt(z)} & The square root \\
\end{booktable}

The four arithmetic operations are written as usual, with \code{+}, \code{-}, \code{*} and
\code{/}.

\subsection{Rectangular and polar form}
\label{c18:sec:polar}
From rectangular to polar form:
\begin{ccode}
const double magnitude = cabs(z);        // |z|
const double angle     = carg(z);        // arg(z) in radians
\end{ccode}
From polar to rectangular form, Euler's formula $z = r e^{j\theta}$ is used:
\begin{ccode}
const double complex z = r * cexp(I * theta);
\end{ccode}

\begin{aside}
\note[Note] \code{printf()} cannot print a complex number directly. The real and imaginary parts
are printed separately:
\begin{ccode}
printf("z = %.2f + %.2fj\n", creal(z), cimag(z));
\end{ccode}
\end{aside}

\section{Impedance and phasors}
\label{c18:sec:impedans}
With complex numbers, impedances are calculated directly in code
(Chapters~\ref{c16:ch}--\ref{c17:ch}). With the angular frequency $\omega = 2\pi f$:
\[
  Z_R = R, \qquad Z_L = j\omega L, \qquad Z_C = \frac{1}{j\omega C} = -\frac{j}{\omega C}
\]
\needspace{7\baselineskip}
\begin{ccode}
const double omega          = 2.0 * PI * f;
const double complex z_r    = r;
const double complex z_l    = I * omega * l;
const double complex z_c    = -I / (omega * c);
const double complex z_tot  = z_r + z_l; // Series connection.
\end{ccode}
A phasor $U = \abs{U} \angle \delta$ is written as a complex number in polar form. Phasors are
then added with ordinary addition:
\begin{ccode}
const double complex u1 = a1 * cexp(I * d1);
const double complex u2 = a2 * cexp(I * d2);
const double complex u  = u1 + u2;
\end{ccode}
The amplitude and phase angle of the result are then given by \code{cabs()} and \code{carg()}
respectively.

\section{Sampling a sinusoidal signal}
\label{c18:sec:sampling}
A sinusoidal voltage is described by
\[
  u(t) = \abs{U} \sin(2\pi f t + \delta),
\]
where $\abs{U}$ is the amplitude in V, $f$ the frequency in Hz, $t$ the time in s and $\delta$ the
phase angle in rad.

In \textbf{sampling}, the signal's value is calculated at evenly spaced times. With the sampling
frequency $f_s$, the number of samples per period is
\[
  N = \frac{f_s}{f},
\]
and the $k$th sample value ($k = 0, 1, \ldots, N-1$) is given by
\[
  u_k = \abs{U} \sin\!\left(\frac{2\pi k}{N} + \delta\right).
\]
For the signal to be recoverable from its samples, the \textbf{sampling theorem} must be
satisfied:
\[
  f_s > 2f
\]
\begin{ccode}
for (int k = 0; k < n; ++k)
{
    const double u_k = amplitude * sin(2.0 * PI * k / n + delta);
    printf("%d,%.6f\n", k, u_k);
}
\end{ccode}

\subsection{Printing for further analysis}
\label{c18:sec:csv}
If the sample values are printed as \textbf{CSV} (comma-separated values), one line per sample,
they can be read into another program to be plotted:
\begin{console}
k,u
0,0.000000
1,3.535534
2,5.000000
\end{console}
In OnlineGDB, select the output in the output box and copy it into a spreadsheet program, where
the signal can be plotted as a chart.

If the program is run on your own computer, the output can instead be saved directly to a file
by redirection in the terminal:
\begin{shell}
./samples > samples.csv
\end{shell}
The file can then be opened in a spreadsheet program, or fed into a simulator that plots the
signal.

\section{Common pitfalls}
\label{c18:sec:fallgropar}
\begin{booktable}{@{}>{\raggedright\arraybackslash}p{10.5em}
    >{\raggedright\arraybackslash}p{10em}L@{}}
  \thead{Pitfall} & \thead{Consequence} & \thead{Remedy} \\
  \midrule
  Forgotten \code{-lm} & \code{undefined reference to 'sqrt'} & Put the flag last when compiling,
    or under \emph{Extra Compiler Flags} in OnlineGDB \\
  Degrees instead of radians & Completely wrong value from \code{sin()}, \code{cos()},
    \code{tan()} & Multiply by $\pi/180$ \\
  Integer division, e.g.\ \mbox{\code{1 / 2}} & Gives \code{0}, not \code{0.5} &
    Write \mbox{\code{1.0 / 2.0}} \\
  \code{atan(y / x)} for polar form & Wrong quadrant & Use \code{atan2(y, x)} \\
  Comparing floating-point numbers with \code{==} & Rarely true & Compare with
    \code{fabs(a - b) < 1e-9} \\
  \code{sqrt()} or \code{log()} with an invalid argument & Gives \code{nan} and \code{-inf}
    respectively & Check the argument first \\
  \code{\%d} for a \code{double} & Undefined behaviour & Use \code{\%f}, for example
    \code{\%.3f} \\
\end{booktable}

\section{Worked examples}
\label{c18:sec:typexempel}

\begin{exempel}[Gain and decibels]
An amplifier has the input voltage $U_{\text{in}} = 0.5\,\text{V}$ and the output voltage
$U_{\text{out}} = 8.0\,\text{V}$. Calculate the gain in dB.
\[
  A_{\text{dB}} = 20\log_{10}\!\left(\frac{8.0}{0.5}\right) = 20\log_{10}(16) \approx
  24.1\,\text{dB}
\]
In code:

\begin{ccode}
const double gain_db = 20.0 * log10(8.0 / 0.5); // 24.082 dB.
\end{ccode}
\end{exempel}

\begin{exempel}[Numerical differentiation]
Differentiate $f(x) = x^2$ at the point $x = 3$ and compare with the analytical derivative.

Analytically, $f'(x) = 2x$, so $f'(3) = 6$. Numerically, with $h = 10^{-6}$:
\[
  f'(3) \approx \frac{(3.000001)^2 - (2.999999)^2}{2 \cdot 10^{-6}} = 6.000000
\]
The deviation is of the order of $10^{-9}$ and is due to the finite precision of floating-point
numbers.
\end{exempel}

\begin{exempel}[Numerical integration]
The current $i(t) = 2t + 4\,\text{A}$ flows for $0 \leq t \leq 2\,\text{s}$. Calculate the charge
$q$.

Analytically (\kapref{14}):
\[
  q = \int_0^2 (2t + 4)\,dt = \Bigl[t^2 + 4t\Bigr]_0^2 = 4 + 8 = 12\,\text{C}
\]
The trapezoidal rule with $n = 4$ and $h = 0.5$:
\[
  q \approx 0.5\left[\frac{4 + 8}{2} + (5 + 6 + 7)\right] = 0.5 \cdot 24 = 12\,\text{C}
\]
The result is exact, because $i(t)$ is a straight line.
\end{exempel}

\begin{exempel}[Impedance of an RL circuit]
A resistor $R = 100\,\Omega$ is connected in series with an inductor $L = 50\,\text{mH}$.
Calculate the circuit's impedance at $f = 50\,\text{Hz}$.
\begin{gather*}
  \omega = 2\pi \cdot 50 \approx 314.16\,\text{rad/s} \\
  Z = R + j\omega L = 100 + j15.71\,\Omega \\
  \abs{Z} = \sqrt{100^2 + 15.71^2} \approx 101.23\,\Omega, \qquad
  \arg(Z) = \arctan\!\left(\frac{15.71}{100}\right) \approx 0.156\,\text{rad} \approx
  8.93^{\circ}
\end{gather*}
In code, the magnitude and the argument are calculated with \code{cabs(z)} and \code{carg(z)}
respectively.
\end{exempel}

\begin{exempel}[Sample values]
A signal $u(t) = 5\sin(2\pi \cdot 50t)\,\text{V}$ is sampled at $f_s = 400\,\text{Hz}$. Calculate
the first three sample values.
\begin{gather*}
  N = \frac{400}{50} = 8 \text{ samples per period} \\
  u_0 = 5\sin(0) = 0\,\text{V} \\
  u_1 = 5\sin\!\left(\frac{2\pi}{8}\right) = 5\sin\!\left(\frac{\pi}{4}\right) \approx
  3.54\,\text{V} \\
  u_2 = 5\sin\!\left(\frac{2\pi \cdot 2}{8}\right) = 5\sin\!\left(\frac{\pi}{2}\right) =
  5.00\,\text{V}
\end{gather*}
The sampling theorem is satisfied, since $400 > 2 \cdot 50$.
\end{exempel}
